\chapter{Sprint 3: Semantic and Parallel Search with Ranking}

\section*{Introduction}
\addcontentsline{toc}{section}{Introduction}

Keyword search alone cannot handle intent-based queries such as ``a football crowd celebrating''.
Sprint 3 adds a \textbf{semantic retrieval} layer to the Videos Search API, runs lexical and
vector searches \textbf{in parallel}, and fuses their rankings with Reciprocal Rank Fusion
(RRF). An optional reranker and a \textbf{position-aware blend} produce the final ranking before
pagination.

%% ─── Sprint Backlog ─────────────────────────────────────────────────────────
\section{Sprint Backlog}

\begin{table}[H]
\centering
\caption{Sprint 3 Backlog}
\label{tab:sprint3_backlog}
\begin{tabularx}{\textwidth}{clXcc}
\toprule
\textbf{ID} & \textbf{Story} & \textbf{Task} & \textbf{Effort (d)} & \textbf{Prio.} \\
\midrule
T3.1 & US-09 & Select embedding model (OpenAI \texttt{text-embedding-3-small})     & 0.5 & M \\
T3.2 & US-09 & Add \texttt{embedding} field (kNN mapping) to the video index       & 1 & M \\
T3.3 & US-09 & Ingestion-side embedding computation and backfill                   & 1 & M \\
T3.4 & US-09 & Semantic search endpoint using OpenSearch \texttt{knn} query        & 1 & M \\
T3.5 & US-10 & Parallel lexical + vector retrieval via \texttt{ThreadPoolExecutor} & 1 & M \\
T3.6 & US-10 & Reciprocal Rank Fusion (RRF) merge with top-rank bonus              & 1 & M \\
T3.7 & US-10 & Low-confidence detector for optional reranking                      & 0.5 & S \\
T3.8 & US-10 & Heuristic-first reranking + OpenAI last-resort reranker            & 1 & S \\
T3.9 & US-10 & Position-aware blend (retrieval \& rerank weights by rank)           & 0.5 & S \\
T3.10 & US-10 & Tunable strategy parameters and \texttt{rerank\_top\_n} cutoff      & 0.5 & C \\
\bottomrule
\end{tabularx}
\end{table}

%% ─── Design ─────────────────────────────────────────────────────────────────
\section{Design}

\subsection{Use Case Diagram for Sprint 3}

Sprint 3 does not introduce a new visible use case. It refines \emph{Search videos} from the
global use-case diagram (Figure~\ref{fig:uc_global}) with semantic and hybrid ranking behaviour.

\subsection{Class Diagram for Sprint 3}

The OpenSearch index gains a vector field. The mapping extract below shows the key additions
to the video document schema.

\begin{lstlisting}[caption={kNN field addition to the OpenSearch video mapping}]
"embedding": {
  "type": "knn_vector",
  "dimension": 1536,
  "method": {
    "name": "hnsw",
    "space_type": "cosinesimil",
    "engine": "lucene"
  }
}
\end{lstlisting}

\subsection{Sequence Diagram: Advanced Video Search}

Figure~\ref{fig:seq_video_search} shows an advanced video-search request traversing the entire
pipeline: the BFF forwards the user query to the Videos Search API, the API fans out parallel
lexical and vector queries, merges the ranked lists with RRF, runs a fast heuristic reranker,
then calls the OpenAI reranker only as a last resort when local confidence is weak, applies the
position-aware blend, and returns a paginated response.

\begin{figure}[H]
\centering
\begin{tikzpicture}[x=1cm, y=1cm, font=\fontsize{7}{8}\selectfont]
%% Participants
\node[sdFE]   (pUser)  at ( 0.0, 0) {User};
\node[sdFE]   (pBFF)   at ( 2.8, 0) {Portal BFF};
\node[sdBE]   (pAPI)   at ( 5.6, 0) {Videos\\Search API};
\node[sdData] (pOS)    at ( 8.4, 0) {OpenSearch};
\node[sdExt]  (pEmb)   at (11.2, 0) {Embedding\\API};
\node[sdExt]  (pRerank)at (14.0, 0) {OpenAI\\reranker};

\foreach \x in {0.0, 2.8, 5.6, 8.4, 11.2, 14.0}
  \draw[dashed, gray!50] (\x,-0.3) -- (\x,-12.0);

\draw[sdmsg] (0.0,-0.85) -- node[above, pos=0.5]{query + filters} (2.8,-0.85);
\draw[sdmsg] (2.8,-1.57) -- node[above, pos=0.5]{POST /api/videos/query} (5.6,-1.57);
\draw[sdmsg] (5.6,-2.29) -- node[above, pos=0.5]{embed(query)} (11.2,-2.29);
\draw[sdret] (11.2,-3.01)-- node[above, pos=0.5]{query\_vector} (5.6,-3.01);

%% Parallel
\node[sdfrag] at (-0.5,-3.65) {par};
\draw[sdmsg] (5.6,-3.73) -- node[above, pos=0.5]{lexical multi\_match} (8.4,-3.73);
\draw[sdmsg] (5.6,-4.45) -- node[above, pos=0.5]{knn vector search} (8.4,-4.45);
\draw[sdret] (8.4,-5.17) -- node[above, pos=0.5]{hits\_lex (ranked)} (5.6,-5.17);
\draw[sdret] (8.4,-5.89) -- node[above, pos=0.5]{hits\_vec (ranked)} (5.6,-5.89);

\draw[sdmsg] (5.6,-6.61) to[out=-30,in=30] node[right, pos=0.5]{RRF fusion + top-rank bonus} (5.6,-7.33);
\draw[sdmsg] (5.6,-8.05) to[out=-30,in=30] node[right, pos=0.5]{heuristic rerank} (5.6,-8.77);

\node[sdfrag] at (-0.5,-9.4) {opt: weak local confidence};
\draw[sdmsg] (5.6,-9.49) -- node[above, pos=0.5]{OpenAI rerank (last resort)} (14.0,-9.49);
\draw[sdret] (14.0,-10.21)-- node[above, pos=0.5]{rerank scores} (5.6,-10.21);

\draw[sdmsg] (5.6,-10.93)-- node[above, pos=0.5]{position-aware blend + sort + paginate} (5.6,-10.93);
\draw[sdret] (5.6,-11.65)-- node[above, pos=0.5]{ranked video results} (0.0,-11.65);

\end{tikzpicture}
\caption{Advanced video search sequence — parallel retrieval and ranking}
\label{fig:seq_video_search}
\end{figure}

%% ─── Implementation ────────────────────────────────────────────────────────
\section{Implementation}

\subsection{Embedding Model and Index Mapping}

Semantic retrieval uses \textbf{OpenAI \texttt{text-embedding-3-small}} (1536-dimensional
embeddings). Each video document now carries an \texttt{embedding} field populated from
\texttt{title + description + tags + transcript (if available)}. The field uses OpenSearch's
\texttt{knn\_vector} type with HNSW indexing and cosine similarity space.

\subsection{Vector Search}

\textbf{Vector search} complements keyword search by matching \textbf{meaning} instead of words.
Each video is represented by its embedding. A query's embedding is compared with the indexed
vectors using cosine similarity through OpenSearch's \texttt{knn} query.

\begin{figure}[H]
\centering
\begin{tikzpicture}[x=1cm,y=1cm,
  box/.style={draw=violet!60!black, fill=violet!8, rounded corners=4pt,
              minimum width=3.4cm, minimum height=0.95cm, align=center, font=\small},
  data/.style={draw=teal!60!black, fill=teal!8, rounded corners=4pt,
               minimum width=3.6cm, minimum height=0.95cm, align=center, font=\small},
  ext/.style={draw=gray!60, fill=gray!8, rounded corners=4pt,
              minimum width=3.0cm, minimum height=0.95cm, align=center, font=\small},
  arr/.style={->, >=Stealth, line width=0.9pt}
]
\node[ext]  (user)  at (-5.8, 0)   {User query};
\node[box]  (embed) at (-1.8, 0)   {OpenAI embedding\\text-embed-3-small};
\node[data] (qvec)  at ( 2.3, 0)   {Query vector\\$q \in \mathbb{R}^{1536}$};
\node[data] (os)    at ( 6.3, 0)   {OpenSearch kNN\\cosine HNSW};
\node[box]  (hits)  at ( 2.3,-2.5) {Top-k neighbours\\by similarity};

\draw[arr] (user) -- (embed);
\draw[arr] (embed) -- (qvec);
\draw[arr] (qvec) -- (os);
\draw[arr] (os.south) |- (hits.east);
\draw[arr] (hits.west) -| node[below, font=\footnotesize, pos=0.25]{semantic results} (user.south);
\end{tikzpicture}
\caption{Vector (semantic) search pipeline}
\label{fig:vector_search}
\end{figure}

Because the comparison is geometric, conceptually similar phrasings produce close vectors even
when they share no tokens. This is what enables natural-language queries to reach relevant videos
that keyword search would have missed.

\subsection{Parallel Retrieval and Ranking Pipeline}

For advanced video discovery, the platform combines lexical and semantic retrieval and then
stabilizes the ranking with lightweight reranking only when needed. The orchestration is
implemented in \texttt{advanced\_search.py}.

\begin{figure}[H]
\centering
\begin{tikzpicture}[
  x=1cm, y=1cm,
  node distance=0.9cm and 1.0cm,
  box/.style={draw=blue!60!black, fill=blue!8, rounded corners=4pt,
              minimum width=3.2cm, minimum height=0.72cm,
              align=center, font=\fontsize{8}{9}\selectfont},
  stage/.style={draw=teal!60!black, fill=teal!8, rounded corners=4pt,
                minimum width=3.2cm, minimum height=0.72cm,
                align=center, font=\fontsize{8}{9}\selectfont},
  dec/.style={draw=orange!70!black, fill=orange!10, diamond,
              aspect=2.2, align=center, inner sep=2pt,
              font=\fontsize{7.5}{8.5}\selectfont},
  arr/.style={->, >=Stealth, thick}
]
\node[box] (norm) {Normalize request\\(query, filters, strategy)};
\node[stage, below=of norm] (subq) {Generate subqueries\\(original + expansion)};

\node[stage, below left=1.2cm and 2.4cm of subq] (lex) {Lexical search\\OpenSearch \texttt{multi\_match}};
\node[stage, below right=1.2cm and 2.4cm of subq] (vec) {Vector search\\OpenSearch \texttt{knn}};

\node[box, below=2.1cm of subq] (rrf) {RRF fusion + top-rank bonus\\candidate pool};
\node[stage, below=of rrf] (heur) {Heuristic rerank\\term overlap};
\node[dec, below=of heur] (gate) {OpenAI\\last resort?};
\node[stage, below left=1.15cm and 2.2cm of gate] (llm) {OpenAI reranker\\score in [0,1]};
\node[box, below right=1.15cm and 2.2cm of gate] (keep) {Keep heuristic\\rerank scores};
\node[box, below=2.4cm of gate] (blend) {Position-aware blend\\retrieval + rerank};
\node[box, below=of blend] (sort) {Sort + paginate\\response items};

\draw[arr] (norm) -- (subq);
\draw[arr] (subq) -- (lex);
\draw[arr] (subq) -- (vec);
\draw[arr] (lex) -- (rrf);
\draw[arr] (vec) -- (rrf);
\draw[arr] (rrf) -- (heur);
\draw[arr] (heur) -- (gate);
\draw[arr] (gate) -- node[left, font=\scriptsize]{yes} (llm);
\draw[arr] (gate) -- node[right, font=\scriptsize]{no} (keep);
\draw[arr] (llm) -- (blend);
\draw[arr] (keep) -- (blend);
\draw[arr] (blend) -- (sort);
\end{tikzpicture}
\caption{Parallel retrieval and ranking pipeline in \texttt{advanced\_search.py}}
\label{fig:advanced_search_ranking}
\end{figure}

\subsubsection{Parallel Retrieval}

Lexical (\texttt{multi\_match}) and vector (\texttt{knn}) searches are dispatched together with a
\texttt{ThreadPoolExecutor}. Two independent candidate pools are collected, each with its own
rank positions. Running them in parallel halves the retrieval latency compared to sequential
execution.

\subsubsection{RRF Fusion}

The two ranked lists are merged using \textbf{Reciprocal Rank Fusion}:
\begin{equation}
\mathrm{RRF}(d)=\sum_{l \in \{\mathrm{lex},\mathrm{vec}\}} \frac{1}{k+\mathrm{rank}_l(d)} + \mathrm{bonus}(\mathrm{rank}_l(d))
\label{eq:rrf}
\end{equation}
with:
\[
\mathrm{bonus}(r)=
\begin{cases}
0.05 & \text{if } r=1\\
0.02 & \text{if } r\in\{2,3\}\\
0 & \text{otherwise}
\end{cases}
\]
where $k$ is \texttt{rrf\_k} (default 60). RRF
uses only the \textit{rank}, not the raw score, so the fusion is robust to score-scale mismatches
between BM25 and cosine similarity.

\subsubsection{Conditional Reranking}

The ranking stage is \textbf{heuristic-first}: token-overlap reranking is computed locally for the
top-$N$ candidates (\texttt{rerank\_top\_n}, default 30). A low-confidence detector then decides
whether to escalate to an OpenAI reranker as a \textbf{last resort} (config:
\texttt{ADVANCED\_SEARCH\_OPENAI\_RERANK\_MODE=last\_resort}). The OpenAI call uses a short timeout
(2 seconds by default) so latency stays bounded; on timeout or API failure the heuristic scores are
kept.

\subsubsection{Position-Aware Blend}

Retrieval trust is higher for the very top results; rerank trust is higher deeper in the list.
We encode this intuition with a rank-dependent weight $w_r$:
\begin{equation}
\mathrm{final}(d)=w_r \cdot \mathrm{retrieval\_norm}(d) + (1-w_r)\cdot \mathrm{rerank}(d)
\label{eq:blend}
\end{equation}
with $w_r=0.8$ for ranks 1--5, $w_r=0.6$ for ranks 6--15, and $w_r=0.4$ for ranks beyond 15.
Default strategy parameters: \texttt{lex\_k=120}, \texttt{vec\_k=120}, \texttt{fuse\_k=150},
\texttt{rerank\_top\_n=30}.

\section*{Conclusion}
\addcontentsline{toc}{section}{Conclusion}

Sprint 3 complements lexical search with an end-to-end semantic pipeline. Parallel retrieval,
rank-based RRF fusion, conditional reranking, and the position-aware blend produce a stable,
low-latency ranker that works well on both precise and intent-driven queries. The next sprint
moves up one abstraction layer and implements the Agentic Helper's RAG and intent classifiers on
top of this search foundation.
